% =====================================================================
%  manuscrito_final.tex -- ARTEFACTO AUTONOMO (version interna no anonima)
%  Proyecto SIMPA -- Equipo AHMRV -- ISR-401 -- UTEQ
%
%  Genero: Research Preview prerregistrado. El estudio NO se ha ejecutado:
%  el documento no contiene resultados ni discusion.
%
%  Este archivo es autocontenido: no depende de preambulo.tex ni de
%  cuerpo_manuscrito.tex. Solo requiere referencias.bib y la clase llncs.
%
%  COMPILACION:
%    pdflatex manuscrito_final
%    bibtex   manuscrito_final
%    pdflatex manuscrito_final
%    pdflatex manuscrito_final
%
%  [DECISION DE AUTORIA PENDIENTE]
%  Los metadatos del proyecto y el deposito declaran seis autores estudiantes.
%  La eventual coautoria del docente supervisor NO esta decidida y por eso no
%  figura. Si se resuelve afirmativamente, hay que anadir la persona al bloque
%  \author y actualizar a la vez manuscrito, repositorio y deposito.
%  Esta nota vive en el fuente y no aparece en el PDF.
% =====================================================================
\documentclass[runningheads]{llncs}
\usepackage[utf8]{inputenc}
\usepackage[T1]{fontenc}
\usepackage[spanish,es-noquoting,es-noshorthands,es-nolists]{babel}
\usepackage{amsmath,amssymb}
\usepackage{url}
\usepackage[hidelinks]{hyperref}
\renewcommand{\keywordname}{\bfseries Palabras clave}
\renewcommand{\refname}{Referencias}

\begin{document}

\title{Calidad de requisitos funcionales: humanos frente a un modelo grande
de lenguaje. Estudio prerregistrado}
\titlerunning{Calidad de requisitos: humanos frente a modelo de lenguaje}

% Sin ORCID: no se declara ninguno hasta disponer de los reales.
\author{
Allan N. Villafuerte Rosero \and
Denisses F. Huilcapi Le\'on \and
Edson N. Rizzo V\'elez \and
Josthyn E. Mac\'ias Herrera \and
Francisco J. Arboleda Yanza \and
Anderson A. Alc\'ivar V\'elez
}
\authorrunning{A. Villafuerte et al.}
\institute{Universidad T\'ecnica Estatal de Quevedo,
Facultad de Ciencias de la Computaci\'on, Quevedo, Ecuador}

\maketitle

% ---------------------------------------------------------------------
% Resumen estructurado en exactamente cuatro parrafos
%   1 contexto y motivacion  2 pregunta o problema
%   3 ideas principales      4 contribucion y limitaciones
% ---------------------------------------------------------------------
\begin{abstract}
\textbf{Contexto y motivaci\'on.} Los modelos grandes de lenguaje se incorporan
con rapidez a la elicitación de requisitos, pero la evidencia emp\'irica sobre la
calidad de los requisitos que producen sigue siendo escasa, y pr\'acticamente
inexistente en dominios agroindustriales y en contextos latinoamericanos.

\textbf{Pregunta o problema.} ?`Difiere la calidad de los requisitos funcionales
elicitados por un equipo humano de la de los generados por un modelo grande de
lenguaje a partir del mismo material fuente, y con qu\'e grado de acuerdo la
valoran personas evaluadoras que desconocen su procedencia?

\textbf{Ideas principales.} Se especifica un estudio comparativo con evaluaci\'on a
ciegas sobre un caso real de ingenier\'ia de requisitos en una unidad productiva de
palma africana. Dos conjuntos de al menos veinticinco requisitos, uno humano y otro
generado por el modelo desde una transcripci\'on de entrevista anonimizada, se
normalizan con identificadores neutros y se punt\'uan de uno a cinco en cinco
dimensiones de calidad por tres personas evaluadoras independientes. El an\'alisis
trata las puntuaciones como ordinales y modela evaluador y requisito como factores
aleatorios cruzados.

\textbf{Contribuci\'on y limitaciones.} La contribuci\'on de esta vista previa es
el dise\~no y su prerregistro p\'ublico, junto con una r\'ubrica operacionalizada
de cinco dimensiones dise\~nada para ser adaptable a otros contextos. El estudio \textbf{a\'un no se
ha ejecutado}: este manuscrito no presenta resultados. Las limitaciones conocidas
son un \'unico material fuente, un \'unico modelo, personas evaluadoras del
entorno acad\'emico y un tama\~no muestral que sit\'ua la potencia por debajo del
objetivo convencional.
\keywords{Ingenier\'ia de requisitos \and Modelos grandes de lenguaje \and
Calidad de requisitos \and Estudio prerregistrado \and Evaluaci\'on a ciegas.}
\end{abstract}

% =====================================================================
\section{Introducci\'on y motivaci\'on}
% =====================================================================

La elicitación de requisitos incide de forma persistente en el resultado de los
proyectos de software \cite{mendez2017}, y los atributos de calidad se definen de
forma heterog\'enea entre estudios \cite{montgomery2022}. Sobre ese terreno se han
incorporado con rapidez los modelos grandes de lenguaje: la elicitación es una de
las fases con menor evidencia acumulada \cite{cheng2026}, valorar la calidad del
artefacto generado exige un procedimiento de evaluaci\'on controlado
\cite{arora2024}, y las pr\'acticas de uso propuestas a\'un no han sido validadas
de forma independiente \cite{vogelsang2025}.

Esta vista previa presenta el dise\~no de un estudio comparativo con evaluaci\'on a
ciegas, prerregistrado antes de recoger puntuaciones, en un contexto poco
representado: un proyecto de ingenier\'ia de requisitos sobre una unidad productiva
agroindustrial real en Ecuador.

% =====================================================================
\section{Trabajo relacionado}
% =====================================================================

La revisi\'on que sustenta este dise\~no es \textbf{dirigida y de car\'acter
narrativo}, no una revisi\'on sistem\'atica: no se conserv\'o un registro auditable
de consultas, y por tanto no se reclama exhaustividad ni reproducibilidad de la
b\'usqueda.

\textbf{Calidad de requisitos.} La ISO/IEC/IEEE 29148:2018 fija los atributos
exigibles a un requisito y a un conjunto \cite{iso29148} y es el marco normativo de
este estudio; Pohl \cite{pohl2010} y Wiegers y Beatty \cite{wiegers2013} los
desarrollan en t\'erminos operativos. La dispersi\'on de definiciones entre estudios
\cite{montgomery2022} motiva que aqu\'i se declare la operacionalizaci\'on de cada
dimensi\'on (Secci\'on~\ref{sec:diseno}). Ferrari et al.\ mostraron que la
ambig\"uedad y el conocimiento t\'acito en entrevistas producen omisiones que el
analista no detecta \cite{ferrari2016}, hallazgo pertinente porque el material
fuente de ambos conjuntos es una transcripci\'on de entrevista. La l\'inea sobre
explicabilidad como requisito no funcional \cite{chazette2021} enmarca los
componentes de inteligencia artificial del sistema objeto del proyecto.

\textbf{Contexto agroindustrial.} La producci\'on reciente sobre palma africana se
concentra en el componente de percepci\'on \cite{omar2024,kipli2023} y no aborda la
especificaci\'on de requisitos del sistema que la incorpora, que es el hueco al que
responde el proyecto en el que se inserta este estudio.

\textbf{M\'etodo.} El dise\~no sigue las gu\'ias de Wohlin et al.\
\cite{wohlin2012} y el reporte la plantilla de Jedlitschka et al.\
\cite{jedlitschka2008}. La pr\'actica de prerregistro se apoya en Nosek et al.\
\cite{nosek2018}.

% =====================================================================
\section{Contexto del estudio}
% =====================================================================

El estudio se inserta en un proyecto de ingenier\'ia de requisitos para un sistema
de gesti\'on del mantenimiento de un cultivo de palma africana, en una unidad
productiva de unas cien hect\'areas en la costa ecuatoriana, designada con un
seud\'onimo por acuerdo de confidencialidad.

La especificaci\'on se elabor\'o conforme a la ISO/IEC/IEEE 29148:2018
\cite{iso29148}, con los atributos de la ISO/IEC 25010:2023 \cite{iso25010} como
marco de los requisitos no funcionales. El cat\'alogo, sometido a inspecci\'on
formal y a control de cambios, comprende \textbf{42 requisitos funcionales, 19 no
funcionales, 10 restricciones de dise\~no y 8 requisitos legales}. La evidencia de
campo procede de entrevistas semiestructuradas, un cuestionario al personal
operativo, documentaci\'on interna y observaci\'on directa, bajo consentimiento
informado y seudonimizada conforme a la legislaci\'on ecuatoriana \cite{lopdp}. El
experimento utiliza \'unicamente la transcripci\'on anonimizada de una entrevista
como material fuente com\'un a los dos conjuntos.

% =====================================================================
\section{Preguntas de investigaci\'on}
% =====================================================================

Marco PICOC: \textbf{poblaci\'on}, requisitos funcionales de un sistema
agroindustrial real; \textbf{intervenci\'on}, su generaci\'on mediante un modelo
grande de lenguaje desde una transcripci\'on anonimizada; \textbf{comparaci\'on},
los elicitados por el equipo humano desde el mismo material; \textbf{resultado}, la
puntuaci\'on en cinco dimensiones de calidad; \textbf{contexto}, un proyecto
acad\'emico de ingenier\'ia de requisitos en Ecuador.

\begin{description}
\item[RQ1.] ?`En qu\'e dimensiones de calidad difieren los requisitos funcionales
elicitados por un equipo humano de los generados por un modelo grande de lenguaje
a partir del mismo material fuente?
\item[RQ2.] ?`Cu\'al es el grado de acuerdo entre personas evaluadoras
independientes al valorar la calidad de requisitos funcionales sin conocer su
procedencia?
\end{description}

Para cada dimensi\'on $d$, la hip\'otesis nula sostiene que la distribuci\'on de
puntuaciones no difiere entre conjuntos. El contraste es \textbf{bilateral}: la
literatura reporta resultados mixtos y postular una direcci\'on sin base
constituir\'ia un sesgo de expectativa.

% =====================================================================
\section{Dise\~no del estudio}
\label{sec:diseno}
% =====================================================================

\subsection{Variables}

La \textbf{variable independiente} es el origen del requisito, nominal
dicot\'omica. Las \textbf{variables dependientes} son cinco puntuaciones ordinales
de 1 a 5, con descriptores por nivel y un ejemplo positivo y otro negativo en cada
uno: \emph{completitud} (presencia de los atributos de la plantilla),
\emph{ausencia de ambig\"uedad} (unicidad de interpretaci\'on),
\emph{verificabilidad} (existencia de un criterio objetivamente comprobable),
\emph{correcci\'on} (fidelidad al contenido de la transcripci\'on) y
\emph{consistencia interna} (ausencia de contradicci\'on con el resto del
conjunto). Se controlan el material fuente y la r\'ubrica, id\'enticos para ambos
conjuntos; la ceguera de las personas evaluadoras respecto del origen; y el orden
de presentaci\'on, aleatorizado con semilla registrada.

\subsection{Materiales y participantes}

Ambos conjuntos comprenden al menos veinticinco requisitos; el humano se extrae del
cat\'alogo seleccionando los trazados a la evidencia de la entrevista fuente.
Se prev\'en tres personas evaluadoras independientes, estudiantes de otro paralelo
de la misma asignatura sin participaci\'on en el proyecto. Se reportar\'a el
n\'umero efectivo, y ese n\'umero constituye una limitaci\'on del estudio.

\subsection{Procedimiento}

Verificada la anonimizaci\'on de la transcripci\'on fuente, el conjunto generado se
obtiene con una consigna literal \'unica, registrando modelo y versi\'on,
temperatura, \emph{top-p}, \emph{top-k}, semilla, fecha y hora, y la respuesta sin
edici\'on \cite{vogelsang2025}. Ambos conjuntos se transcriben a una plantilla
com\'un con identificadores neutros, sin marca de procedencia, y se mezclan en orden
aleatorio con semilla registrada; la correspondencia entre identificador y origen
se conserva aparte y no se entrega a las evaluadoras. Cada una punt\'ua todos los
requisitos de forma independiente. La procedencia se revela solo al completarse la
recolecci\'on.

% =====================================================================
\section{Plan de an\'alisis}
\label{sec:analisis}
% =====================================================================

\subsection{Desviaci\'on declarada respecto del prerregistro}

El protocolo prerregistrado prev\'e contrastes \textbf{apareados} sobre unos
veinticinco pares. Al revisar el dise\~no antes de recoger puntuaci\'on alguna se
detect\'o que \textbf{no existe correspondencia uno a uno} entre cada requisito
humano y uno generado: ambos conjuntos proceden del mismo material pero no
describen necesariamente las mismas funciones, de modo que el emparejamiento carece
de fundamento y el contraste apareado ser\'ia inv\'alido.

Se declara por tanto la siguiente desviaci\'on, que se someter\'a como
\textbf{adenda prospectiva al registro} \emph{antes} de recoger ninguna
puntuaci\'on:

\begin{enumerate}
\item Los conjuntos se tratan como \textbf{grupos independientes}, no apareados.
\item El an\'alisis principal emplea un \textbf{modelo mixto de v\'inculo
acumulativo} para respuesta ordinal \cite{agresti2010}, que respeta la naturaleza
ordinal de la escala y la estructura de evaluaciones repetidas, que un contraste
sobre medias no captura. Se ajustar\'a \textbf{un modelo por cada una de las cinco
dimensiones}, con la especificaci\'on
$\text{puntuaci\'on}_{1\text{--}5} \sim \text{origen} + (1 \mid \text{evaluador}) + (1 \mid \text{requisito})$,
donde el origen es efecto fijo y evaluador y requisito son interceptos aleatorios
cruzados. Se emplear\'a enlace log\'istico de probabilidades proporcionales, salvo
que antes de ejecutar exista una raz\'on metodol\'ogica preespecificada para otro
enlace. La implementaci\'on prevista es el paquete \texttt{ordinal} de R, cuya
versi\'on exacta se registrar\'a en el momento de la ejecuci\'on.
\item Si el modelo no converge, el an\'alisis de reserva es la prueba
\textbf{U de Mann--Whitney} sobre la mediana por requisito entre evaluadoras, con
$\delta$ de Cliff \cite{cliff1993} e intervalo de confianza al 95\,\% como tama\~no
del efecto.
\end{enumerate}

Ninguna de estas decisiones se tomar\'a en funci\'on de los datos observados.

\subsection{Contrastes y acuerdo entre evaluadoras}

Se contrastan cinco dimensiones sobre el mismo conjunto de datos; la
correcci\'on secuencial de Holm \cite{holm1979} se aplica a los \textbf{cinco
contrastes del efecto de origen, uno por dimensi\'on}, y se reportan los valores
$p$ crudos y corregidos, con $\alpha = 0{,}05$. Ninguna de estas decisiones se
revisar\'a tras observar las puntuaciones.

Para RQ2, el acuerdo se estima con el \textbf{$\alpha$ de Krippendorff con
m\'etrica ordinal} \cite{hayes2007} como medida principal, y con \textbf{$\kappa$
ponderado cuadr\'atico} \cite{cohen1968} por pares como medida secundaria. Se
descarta el $\kappa$ nominal de Cohen o de Fleiss previsto en el registro
original: sobre una escala ordinal de 1 a 5 penaliza por igual un desacuerdo de un
punto y uno de cuatro. Este cambio forma parte de la adenda descrita arriba.

\subsection{Potencia}

Como \textbf{aproximaci\'on de sensibilidad}, y no como potencia del modelo
ordinal, se calcul\'o la potencia de un contraste convencional de dos grupos
independientes con $\alpha = 0{,}05$ bilateral y $d = 0{,}5$: veinticinco
requisitos por grupo rinden aproximadamente \textbf{0,41}, y alcanzar 0,80
requerir\'ia unos 64 por grupo. Una evaluaci\'on de potencia espec\'ifica para el
modelo mixto ordinal exigir\'ia simulaci\'on, una vez preespecificadas las
distribuciones y las componentes de varianza necesarias. La limitaci\'on se declara
antes de ejecutar y no se corregir\'a ampliando la muestra a posteriori: la
interpretaci\'on primar\'a el tama\~no del efecto con su intervalo de confianza
sobre la significaci\'on estad\'istica.

Todas las tablas y figuras del futuro reporte se producir\'an mediante
\emph{scripts} versionados; no se admitir\'an cifras calculadas manualmente.

% =====================================================================
\section{Amenazas a la validez}
% =====================================================================

Se identificaron antes de la ejecuci\'on, seg\'un las categor\'ias de Wohlin et
al.\ \cite{wohlin2012}.

\textbf{Constructo.} La r\'ubrica puede no capturar la calidad de un requisito en su
totalidad; las dimensiones proceden de la literatura \cite{iso29148,montgomery2022} y
la limitaci\'on se declara. La escala ordinal comprime juicios heterog\'eneos, lo que
se mitiga con descriptores y ejemplos por nivel.

\textbf{Interna.} El conjunto humano lo elabor\'o el mismo equipo que ejecuta el
estudio; se mitiga con evaluaci\'on a ciegas, identificadores neutros y orden
aleatorizado con semilla registrada.

\textbf{Externa.} Un \'unico material fuente, dominio y modelo; no se generalizar\'a
m\'as all\'a del caso. Las evaluadoras son estudiantes y no profesionales: se declara
el perfil real y no se atribuir\'a a su juicio la condici\'on de experto.

\textbf{Conclusi\'on.} La potencia queda por debajo de 0,80, por lo que se reporta
el tama\~no del efecto con intervalo de confianza y no solo el valor $p$. Un acuerdo
bajo entre evaluadoras se reportar\'a igualmente y se discutir\'a como hallazgo.

% =====================================================================
\section{Estado y contribuci\'on esperada}
\label{sec:estado}
% =====================================================================

El protocolo est\'a \textbf{dise\~nado y prerregistrado p\'ublicamente}, con marca
temporal anterior a la ejecuci\'on del experimento y a la recolecci\'on de
puntuaciones. El material fuente procede de trabajo de campo previo, ya existente
en el momento del registro; el prerregistro no lo antecede ni lo pretende.
\textbf{La ejecuci\'on est\'a pendiente}, y por ello este manuscrito no contiene
resultados ni discusi\'on. Un registro p\'ublico no equivale a un \emph{Registered
Report} con aceptaci\'on en fase~1: este trabajo es una vista previa de
investigaci\'on prerregistrada, sin revisi\'on previa del dise\~no.

El equipo se compromete a reportar el resultado con independencia de si confirma o
refuta la hip\'otesis nula, a no modificar el plan seg\'un los datos observados y a
declarar el n\'umero efectivo de evaluadoras y de requisitos. La contribuci\'on
esperada es doble: evidencia emp\'irica en un dominio agroindustrial poco
representado, y una r\'ubrica operacionalizada de cinco dimensiones con sus
descriptores.

% =====================================================================
\section*{Declaraci\'on de disponibilidad de datos}
% =====================================================================


El prerregistro del protocolo est\'a disponible en Open Science Framework
(\url{https://osf.io/4z35d/}). El repositorio del proyecto, con la especificaci\'on
de requisitos y la matriz de trazabilidad, es p\'ublico
(\url{https://github.com/AlanNVR/SIMPA_ISR401}).

Como material de apoyo se ha depositado un conjunto con estad\'istica descriptiva
agregada del cuestionario aplicado al personal operativo
(\href{https://doi.org/10.5281/zenodo.22236500}{10.5281/zenodo.22236500}).
\textbf{Ese dep\'osito no contiene los datos del experimento aqu\'i descrito}, que
todav\'ia no se ha ejecutado.


Se publicar\'an \textbf{cuando se ejecute el estudio}: la consigna literal y los
par\'ametros de generaci\'on, los dos conjuntos normalizados, la r\'ubrica completa
con sus descriptores y los \emph{scripts} de an\'alisis. A la fecha de este
manuscrito el repositorio contiene la plantilla de registro de consignas, sin
ejecuciones registradas, y los instrumentos de forma parcial. El \emph{commit} de
ejecuci\'on se fijar\'a y declarar\'a en su momento: ninguna etiqueta actual
identifica una versi\'on evaluada del experimento.

Las puntuaciones individuales de cada evaluadora \textbf{solo se publicar\'an si su
consentimiento espec\'ifico autoriza ese nivel de detalle}; en caso contrario, solo
datos agregados que no permitan atribuir juicios a personas concretas. El dise\~no
de publicaci\'on sigue los principios FAIR \cite{wilkinson2016}.

% =====================================================================
\section*{Uso de tecnolog\'ias asistidas por inteligencia artificial}
% =====================================================================

\textbf{Como objeto de estudio.} La generaci\'on de uno de los conjuntos mediante un
modelo grande de lenguaje es la intervenci\'on experimental; el modelo y su
versi\'on, los par\'ametros, la consigna literal y la respuesta sin editar se
registrar\'an y publicar\'an.

\textbf{Como apoyo a la redacci\'on.} Las personas autoras emplearon el asistente
Claude (Anthropic) en la redacci\'on y reestructuraci\'on de este manuscrito, en la
revisi\'on de coherencia interna del texto y en la elaboraci\'on de un borrador de
la lista de referencias, incluida la sugerencia de obras metodol\'ogicas
can\'onicas. \textbf{Antes del env\'io, las referencias citadas ser\'an contrastadas contra sus
fuentes bibliogr\'aficas o editoriales originales y sus identificadores
persistentes cuando est\'en disponibles.} El asistente no se emple\'o para
generar resultados ni para redactar interpretaciones de datos, que no existen. La
responsabilidad sobre el contenido, las decisiones metodol\'ogicas y las
afirmaciones sostenidas es \'integramente de las personas autoras.


% =====================================================================
\section*{Agradecimientos y conflicto de intereses}
% =====================================================================

Las personas autoras agradecen a la unidad productiva participante y a su personal
el tiempo dedicado a las entrevistas y el acceso a la documentaci\'on operativa. Una
de las personas autoras mantiene una relaci\'on de parentesco con el representante de
la organizaci\'on participante; la circunstancia se declar\'o ante la instancia
acad\'emica antes del trabajo de campo y se adoptaron medidas de mitigaci\'on
documentadas. No se declaran otros conflictos de intereses.


\bibliographystyle{splncs04}
\bibliography{referencias}

\end{document}
